% !TeX root = ../../Book3.tex
\section{椭圆函数的基本结构}
\label{b3:sec:2202}

平行四边形的对边积分相消，把局部极点资料转成全局限制。以下各积分均选取无零极点的边界，并把零极点按模格的重数计数。

\subsection{零点与极点}
\label{b3:subsec:220201}

\begin{proposition}\label{b3:prop:22-zero-pole-count}
非恒零椭圆函数在一个基本区域内的零点总重数等于极点总阶数。
\end{proposition}
\begin{proof}
对数导数 \(f'/f\) 仍双周期，对边方向相反，所以边界积分为零。辐角原理将积分除以 \(2\pi\mathrm{i}\) 识别为零点数减极点数。商空间紧致保证两类点都有限；等式不包含恒零函数。
\end{proof}
对任意有限值 \(a\)，同一结论应用于 \(f-a\)，说明取值 \(a\) 的原像按重数计数等于极点总阶数。这个整数对 \(a\) 不变，给出椭圆函数作为覆盖映射的次数。

\subsection{留数和}
\label{b3:subsec:220202}

\begin{proposition}\label{b3:prop:22-residue-sum}
椭圆函数在一个基本区域内的留数之和为零。
\end{proposition}
\begin{proof}
将上一证明的被积函数改成 \(f\)，用留数定理而不是辐角原理；对边仍因周期性相消。
\end{proof}
因此一个基本区域中不能恰有一个一阶极点，因为一阶极点的留数非零。但一个二阶极点可以出现，其留数可以为零。下一节的 \(\wp\) 函数正是次数二的最简单实例。两个一阶极点的留数必须相反；这个条件将允许用两个平移的对数导数相减来构造它。

\subsection{除数关系}
\label{b3:subsec:220203}

\begin{proposition}\label{b3:prop:22-divisor-necessary}
设非恒零椭圆函数的零点代表为 \(a_j\)，极点代表为 \(b_k\)，均按重数重复。则
\(\sum_j a_j-\sum_k b_k\in\Lambda\)。
\end{proposition}
\begin{proof}
令 \(h=f'/f\)。留数定理给差值等于 \((2\pi\mathrm{i})^{-1}\int_{\partial P}zh(z)\,\mathrm{d}z\)。将对边配对，得到两个基周期乘以沿另两条边的 \(h\) 积分，并附相应正负号。每一条这样的边，其端点处 \(f\) 值相同；沿边连续选对数，终点与起点对数相差 \(2\pi\mathrm{i}\) 的整数倍。因此所得差是 \(\omega_1,\omega_2\) 的整数线性组合。换代表只改变格元素，不影响结论。
\end{proof}
这一条件比零极点数相等更强。例如指定单零点 \(a\) 和单极点 \(b\) 会要求 \(a-b\in\Lambda\)，这时零极点在商空间上相同并相互抵消，所以次数一不可能出现。给定除数的存在性则还需要构造函数，不能从必要条件倒推；下文先建立足以进行这种构造的 \(\wp\) 及其积分。

\subsection{椭圆函数的次数}
\label{b3:subsec:220204}

若非常值椭圆函数在基本区域有总阶数 \(d\) 的极点，就称 \(d\) 为它相对于 \(\Lambda\) 的\term[tuoyuan-hanshu-cishu]{次数}[degree of an elliptic function]。局部展开 \(f(z)-a=c(z-z_0)^m+\cdots\) 中的 \(m\) 是该纤维点的重数；极点处改用 \(1/f\) 计算。

\ProseParagraphBreak

在非分歧值附近，各个局部逆分支连续存在，故原像个数恰为 \(d\)。分歧值上不同逆分支可能合并，但总重数不变。这样，辐角原理的整数不只是记账工具，而是刻画了映射的每条纤维。次数二的偶函数预示了一对原像 \(z,-z\)，它们只在半周期点合并。
